\section{A bounded lattice on the test sphere}
\label{sec:lattice}

Set
\[
 F_r(X)=\Witt_r(\Pcat(X)),\qquad
 E_r(X)=\Witt_r(\T(X)).
\]
The inclusion of categories gives a natural comparison
$F_r(X)\to E_r(X)$.  In this section spheres and their pole complements
have their usual semialgebraic models.  Rationalization always means
tensoring these abelian groups with $\Q$; the sheaf coefficients remain
$\R$.  Fix an odd integer $d\geq3$.  We will prove that the rational
image of $E_d(S^d)$ lies in a lattice with one denominator $L_d$.
The bound may depend on $d$, but will be independent of every action,
test map, and character partition used later.

\subsection{Point groups and finite exponent comparison}

\begin{lemma}[The point calculation]
\label{lat:point-groups}
For every integer $r$, the comparison at a point is an isomorphism, and
\[
 F_r(\pt)=E_r(\pt)=
 \begin{cases}
  \Z,&r\equiv0\pmod4,\\
  0,&r\not\equiv0\pmod4.
 \end{cases}
\]
In degree zero the indicated isomorphism is signature.  On $S^0$ both
theories are the direct sum of two copies of the point theory, and the
map to the point is addition.
\end{lemma}

\begin{proof}
A generator over a point is the finite dimensional bounded cohomology
complex of a compact finite polyhedron.  Conversely, the point constant
and its shifts generate every bounded complex of finite dimensional
real vector spaces: split a complex into its cohomology and a
contractible complex.  Both point categories are therefore precisely
this derived category.

For a degree $r$ form, the cohomology pairing pairs $H^i$ perfectly with
$H^{r-i}$.  The complex splitting identifies the form with the
orthogonal sum of these cohomology pairings.  Terms with $i\ne r-i$
occur in hyperbolic pairs: choose one member as a lagrangian, the other
being its shifted dual.  If $r$ is odd there is no remaining term.
If $r=2k$, the remaining form is on $H^k$, with symmetry sign
$(-1)^{k^2}=(-1)^k$ from tensor symmetry.  For odd $k$ it is alternating;
a symplectic basis exhibits it as hyperbolic.  For even $k$ it is a
real symmetric form.  Diagonalization and cancellation of a positive
and a negative line give its class as its signature times the
positive line.

Signature respects all neutral relations.  Indeed, in a lagrangian
triangle $L\to A\to\D_{\pt,r}L\to L[1]$, at the middle degree $k$ the
second map is the adjoint of the first for the middle pairing.  Thus
exactness identifies the image of $H^k(L)$ with its own perpendicular
in $H^k(A)$.  Such a symmetric form has signature zero.  The positive
line has signature one, proving the claimed infinite cyclic group.
Finally, sheaves and their dualities on the two-point space split by
the two open and closed components, giving the last assertion.
\end{proof}

We record the elementary algebra that will keep the denominator bound
uniform.  The class of abelian groups annihilated by some finite power
of $2$ is closed under subgroups, quotients, and finite extensions:
an extension of groups annihilated by $2^a$ and $2^b$ is annihilated by
$2^{a+b}$.  The exponent is allowed to depend on the group, but not on
an element of it.

\begin{lemma}[Finite diagram chases]
\label{lat:finite-chases}
In a commutative diagram of exact rows
\[
\begin{array}{ccccccccc}
 A_1&\longrightarrow&A_2&\longrightarrow&A_3&\longrightarrow&A_4&\longrightarrow&A_5\\
 {\scriptstyle f_1}\downarrow&&{\scriptstyle f_2}\downarrow&&
 {\scriptstyle f_3}\downarrow&&{\scriptstyle f_4}\downarrow&&
 {\scriptstyle f_5}\downarrow\\
 B_1&\longrightarrow&B_2&\longrightarrow&B_3&\longrightarrow&B_4&\longrightarrow&B_5,
\end{array}
\]
suppose the kernels and cokernels of $f_1,f_2,f_4,f_5$ are annihilated
by $q$.  Then the kernel and cokernel of $f_3$ are annihilated by $q^3$.

Suppose also that a commutative square is written as
\[
\begin{array}{ccc}
 A&\xrightarrow{i}&A'\\
 {\scriptstyle f}\downarrow&&\downarrow{\scriptstyle g}\\
 B&\xrightarrow{j}&B'.
\end{array}
\]
If both horizontal kernels and cokernels are annihilated by $2$, and
the kernel and cokernel of $f$ are annihilated by $q$, then those of
$g$ are annihilated by $4q$.
\end{lemma}

\begin{proof}
Denote the arrows in the first rows by $a_i$ and $b_i$.
If $x\in\ker f_3$, then $a_3(qx)=0$, so choose $v\in A_2$ with
$a_2v=qx$.  Since $b_2f_2v=0$, choose $w\in B_1$ with
$b_1w=f_2v$.  Lift $qw$ to $z\in A_1$, so that $f_1z=qw$.
Now $qv-a_1z\in\ker f_2$, hence $q(qv-a_1z)=0$.  Applying $a_2$
gives $q^3x=0$.

For $y\in B_3$, choose $v\in A_4$ with $f_4v=q b_3y$.
Then $a_4v\in\ker f_5$, so choose $w\in A_3$ with $a_3w=qv$.
The element $q^2y-f_3w$ lies in $\ker b_3$; write it as $b_2z$
with $z\in B_2$.  Choose $t\in A_2$ with $f_2t=qz$.  The equality
\[
 f_3(qw+a_2t)=qf_3w+qb_2z=q^3y
\]
proves the cokernel assertion.

In the square, if $g(x')=0$, lift $2x'$ to $x\in A$.  Then
$jf(x)=0$, so $f(2x)=0$ and $2qx=0$.  Applying $i$ gives
$4qx'=0$.  For $y'\in B'$, lift $2y'$ to $y\in B$ and then lift
$qy$ to $x\in A$.  This gives $2qy'=g(ix)$, which even annihilates
the cokernel by $2q$.
\end{proof}

\begin{samepage}
\begin{theorem}[Bounded comparison]
\label{thm:bounded-comparison}
For every $j\ge0$ there is a nonnegative integer $e_j$, independent
of $r$, such that the kernel and cokernel of
\[
 F_r(S^j)\longrightarrow E_r(S^j)
\]
are annihilated by $2^{e_j}$ for every integer $r$.
\end{theorem}
\end{samepage}

\begin{proof}
The assertion for $j=0$ follows from \Cref{lat:point-groups}, with
$e_0=0$.  Suppose it holds for $j-1$.  Put $X=S^j$, let $U,V$
be the complements of two distinct poles, and put $I=U\cap V$.
The spaces $U,V$ are semialgebraically contractible, and $I$ is
semialgebraically homotopy equivalent to $S^{j-1}$ (for $j=1$ this
has two components).  By \Cref{thm:homotopy}, the comparisons on
$U,V$ are isomorphisms and those on $I$ have the induction bound
in every degree.  These homotopies need not be proper: the
pushforward functors of \Cref{prop:test-pushforward} use compact
supports.

For $\Bcat=\Pcat$ or $\T$ write temporarily
\[
 Q_r^{\Bcat}(Z,A)=\Witt_r\bigl(\Bcat(Z)/\Bcat(A)\bigr)
\]
when $A$ is open in $Z$.  Naturality of
\Cref{thm:witt-localization} gives two exact rows of the form
\[
 \Witt_r(\Bcat(I))\longrightarrow\Witt_r(\Bcat(V))
 \longrightarrow Q_r^{\Bcat}(V,I)
 \longrightarrow\Witt_{r-1}(\Bcat(I))
 \longrightarrow\Witt_{r-1}(\Bcat(V)).
\]
Their four outer comparisons have the uniform induction bound.
The first chase in \Cref{lat:finite-chases} therefore bounds the
comparison on $Q_r(V,I)$ in all degrees.

By \Cref{prop:open-cover,thm:dense-doubling}, the two horizontal maps
in the natural square
\[
\begin{array}{ccc}
 Q_r^{\Pcat}(V,I)&\longrightarrow&Q_r^{\Pcat}(X,U)\\
 \downarrow&&\downarrow\\
 Q_r^{\T}(V,I)&\longrightarrow&Q_r^{\T}(X,U)
\end{array}
\]
have kernels and cokernels annihilated by $2$.  The square chase
therefore gives a uniform finite exponent for the right-hand
comparison.  Finally use the two localization rows
\[
 Q_{r+1}^{\Bcat}(X,U)\longrightarrow\Witt_r(\Bcat(U))
 \longrightarrow\Witt_r(\Bcat(X))
 \longrightarrow Q_r^{\Bcat}(X,U)
 \longrightarrow\Witt_{r-1}(\Bcat(U)).
\]
A further five-term chase gives the desired bound on the middle
comparison.  Each step used only a finite exact diagram, and the
same bounds work for every $r$.
\end{proof}

\begin{remark}
\label{lat:exponent-bookkeeping}
The preceding proof can be made quantitative without optimizing any
bound: if the induction exponent is $e$, its three steps give
$3e$, then $3e+2$, then $9e+6$.  Thus one may take
$e_0=0$ and $e_j=9e_{j-1}+6$, or equivalently
$e_j=3(9^j-1)/4$.  This bound depends on the sphere dimension $j$,
but works for every degree and every class in the comparison groups.
Only existence of such a fixed exponent will be used.
A comparison merely after inverting $2$ would not
suffice: the inclusion $\Z\subset\Z[1/2]$ has an unbounded
$2$-primary torsion cokernel.
\end{remark}

\subsection{The rational sphere group and its integral generator}

\begin{proposition}[Rational sphere calculation]
\label{prop:sphere-rank}
For either $G_r=F_r$ or $G_r=E_r$, and for $j\ge0$,
\[
 G_r(S^j)\otimes\Q\simeq
 \bigl(G_r(\pt)\otimes\Q\bigr)
 \oplus\bigl(G_{r-j}(\pt)\otimes\Q\bigr).
\]
In particular, for the fixed odd integer $d$, both
$F_d(S^d)\otimes\Q$ and $E_d(S^d)\otimes\Q$ are
one-dimensional rational vector spaces.
\end{proposition}

\begin{proof}
Write $\widetilde G_r(Z)=\ker(G_r(Z)\to G_r(\pt))$ for the
kernel of projection to a point.  A point inclusion splits this
projection whenever $Z$ is nonempty.
Write $\Bcat$ for the category giving $G$, and use the notation
$Q_r^{\Bcat}$ from the preceding proof.  For $j\ge1$, use its
pole cover $X=U\cup V$ and $I=U\cap V$.
Since $G_r(U)\to G_r(\pt)$ is an isomorphism, the map
$G_r(U)\to G_r(X)$ is split injective in every degree.  Localization
therefore identifies
\[
 Q_r^{\Bcat}(X,U)\simeq\widetilde G_r(X).
\]
More explicitly, its kernel in $G_r(X)$ is exactly the point
summand, and its next boundary is zero because
$G_{r-1}(U)\to G_{r-1}(X)$ is injective.

Likewise $G_r(I)\to G_r(V)$ is split surjective: project to the
point, use contractibility of $V$, and choose a point of $I$ as a
section.  The other localization sequence consequently identifies
its boundary as an isomorphism
\[
 Q_r^{\Bcat}(V,I)\simeq\widetilde G_{r-1}(I).
\]
The dense map between these two relative groups is an isomorphism
after tensoring with $\Q$ by \Cref{thm:dense-doubling}.  Thus
\[
 \widetilde G_r(S^j)\otimes\Q
 \simeq\widetilde G_{r-1}(S^{j-1})\otimes\Q.
\]
For $S^0$ the reduced group is the kernel of addition of two point
groups, hence one point group.  Induction and the split point
summand prove the formula.  At $(r,j)=(d,d)$ the two summands are
$0$ and $\Q$, respectively, because $d$ is odd and
\Cref{lat:point-groups} gives $G_d(\pt)=0$ and $G_0(\pt)=\Z$.
\end{proof}

Give $S^d$ an orientation.  Its constant real sheaf, with its
orientation pairing into $\omega_{S^d}[-d]$, determines
\[
 u_{F,d}\in F_d(S^d),\qquad u_d\in E_d(S^d).
\]
The comparison takes $u_{F,d}$ to $u_d$.  Write
$\overline{u}_{F,d}=u_{F,d}\otimes1$ and
$\overline{u}_d=u_d\otimes1$ for their rational images.

\begin{proposition}[Signature on a finite semialgebraic relation]
\label{prop:generic-signature}
Every finite collection of objects of $\Pcat(S^d)$ has finite
dimensional locally constant cohomology on a common dense open
subset of $S^d$.  On a sufficiently small oriented ball there, a
degree-$d$ form has an integral signature on its degree-zero
cohomology.  For any finite
Witt relation one can choose such a ball simultaneously for all
objects and lagrangian diagrams in the relation; the signatures
then respect that relation.

Consequently $u_{F,d}$ is nontorsion and
\[
 \begin{gathered}
 F_d(S^d)\otimes\Q=\Q\overline{u}_{F,d},\\
 \im\bigl(F_d(S^d)\to F_d(S^d)\otimes\Q\bigr)=\Z\overline{u}_{F,d}.
 \end{gathered}
\]
\end{proposition}

\begin{proof}
For a continuous semialgebraic map from a compact polyhedron,
Hardt's semialgebraic triviality theorem gives a finite
semialgebraic partition of the target over whose pieces the map
is trivial \cite[main theorem]{Hardt1980}; see also the precise
finite-partition formulation in \cite[Theorem~4.1]{Coste2000}. Refining
this partition into cells,
the union of the top dimensional cells is open and dense in the
sphere.  Each fiber is a compact semialgebraic set, so has the
finite cohomology of a finite triangulation; these semialgebraic
stratification and triangulation facts are recalled in
\cite[\S2.1, pp.~2481--2483]{HardtLambrechtsTurchinVolic2011}.
Proper base change
and the trivializations show that a generator has finite
dimensional locally constant cohomology on that open set.

This property passes to the finite constructions defining
$\Pcat$.  Shifts and summands cause no difficulty.  For a cone,
the cohomology long exact sequence expresses each cohomology
sheaf as an extension of a kernel and a cokernel of morphisms of
finite local systems.  Those kernels and cokernels are local
systems, and such extensions are locally constant: on a small
ball, positive extension groups between finite constant sheaves
vanish.  Indeed, for finite vector spaces $V,W$ and a ball $B$,
\[
 \operatorname{Ext}^k_B(\const{V}{B},\const{W}{B})
 =H^k\bigl(B;\const{\Hom_{\R}(V,W)}{B}\bigr)=0\qquad(k>0).
\]
Taking the finite intersection of the dense open subsets needed
for a finite list of generators proves the first assertion for
every finite list of objects.  No condition is imposed on the
derived morphisms in their constructions.

Choose a ball $B$ in that common open subset.  Every local system
in question is constant on $B$, and the same extension vanishing
splits the successive truncation triangles of each bounded
complex.  Hence its restriction is isomorphic to the constant
complex associated to its finite dimensional cohomology complex.
In particular the constant-complex functor identifies the full
subcategory of bounded complexes on $B$ with finite locally
constant cohomology with the bounded derived category of finite
dimensional vector spaces.  The orientation
identification $\omega_B\simeq\const{\R}{B}[d]$ identifies
$\D_{B,d}$ on this subcategory with ordinary derived linear
duality.  This identification is compatible with evaluation and
biduality: on a finite constant complex it is the natural map
\[
 \const{\operatorname{RHom}_{\R}(C,\R)}{B}
 \simeq
 \operatorname{R\mathcal Hom}_B(\const{C}{B},\const{\R}{B}),
\]
and the assertion passes through its finite shifts and sums.
The orientation form of the constant line becomes the positive
form $xy$.  This degree-zero local description holds in both odd
congruence classes of $d$ modulo four.

A nonsingular degree-$d$ form on $A$ thus gives an ordinary
nonsingular symmetric form $b$ on $H^0(A)_x$, for $x\in B$.
Its signature is an integer, additive under orthogonal sums and
invariant under isometry.  To check neutrality, include a
lagrangian triangle in the finite list before choosing $B$.
Writing $\ell:H^0(L)_x\to H^0(A)_x$, its middle exact sequence is
\[
 H^0(L)_x\xrightarrow{\ell}H^0(A)_x
 \xrightarrow{\ell^*b}\bigl(H^0(L)_x\bigr)^*.
\]
The last map can have the harmless common sign of the triangle
convention.  Exactness gives
$\im\ell=\ker(\ell^*b)=(\im\ell)^\perp$.
This is an ordinary lagrangian; the middle signature is zero.
Since each equality in a Witt group is witnessed by a finite
combination of isometries, orthogonal sums, and neutral
relations, this proves the assertion about any fixed relation.

The class $u_{F,d}$ has signature one on every such oriented ball.
If a nonzero multiple of it vanished, a finite relation witnessing
that equality would contradict its signature.  Thus $u_{F,d}$ is
nontorsion and spans the rational line of \Cref{prop:sphere-rank}.
For $x\in F_d(S^d)$ write $x\otimes1=(a/b)\overline{u}_{F,d}$, where $a\in\Z$
and $b>0$.  The difference $bx-au_{F,d}$ is torsion, so for some
integer $\nu>0$ there is an actual equality
$\nu bx=\nu au_{F,d}$ in the Witt group.
Choose a common ball for this finite equality and put $\sigma$
for the signature of a representative of $x$ there.  It gives
$\nu b\sigma=\nu a$, hence $a/b=\sigma\in\Z$.  All integral multiples
of $u_{F,d}$ occur, proving the asserted equality of images.
\end{proof}

\begin{theorem}[The fixed sphere lattice]
\label{thm:sphere-lattice}
Let $d\geq3$ be odd, with orientation class $u_d$ as above.
There is a positive power of $2$,
denoted by $L_d$, such that
\begin{equation}
 \label{lat:lattice}
 \begin{gathered}
 E_d(S^d)\otimes\Q=\Q\overline{u}_d,\\
 \im\bigl(E_d(S^d)\to E_d(S^d)\otimes\Q\bigr)
 \subseteq L_d^{-1}\Z\overline{u}_d.
 \end{gathered}
\end{equation}
One may take $L_d=2^{e_d}$.  This denominator is fixed by the sphere
dimension and comparison.  It is independent of every later action,
test map, character partition, prime $p$, and exponent $k$.
\end{theorem}

\begin{proof}
By \Cref{thm:bounded-comparison} the comparison is a rational
isomorphism, so \Cref{prop:generic-signature} identifies its
nonzero rational orientation class with $\overline{u}_d$.  Take $L_d=2^{e_d}$,
which annihilates the comparison cokernel in degree $d$.
For every $x\in E_d(S^d)$ there is $y\in F_d(S^d)$ with
$L_dx$ equal to its
image.  The rational image of $y$ is an integral multiple of $\overline{u}_{F,d}$ by
\Cref{prop:generic-signature}, so
$x\otimes1\in L_d^{-1}\Z\overline{u}_d$.
\end{proof}

\subsection{Degree and integral localization}

Two further properties will connect the lattice to the action: degree
compares sphere tests, and integral localization lifts the character
classes from a quotient category.

\begin{proposition}[Action of degree]
\label{prop:degree-action}
Let $d\geq3$ be odd. If $f:S^d\to S^d$ is continuous and has
degree $s\in\Z$, then
\[
 f_*\overline{u}_d=s\overline{u}_d\quad\text{in }E_d(S^d)\otimes\Q.
\]
Consequently $f_*$ acts on this entire rational group by
multiplication by $s$.  More generally, if $N$ is a connected
closed oriented triangulable $d$-manifold and $f:N\to S^d$
has degree $s$, its pushed orientation form has rational class
$s\overline{u}_d$.
\end{proposition}

\begin{proof}
Choose a finite polyhedral model for the source and a finite
semialgebraic triangulation of the target sphere.  The latter
identifies its polyhedral model with $S^d$ semialgebraically.
On these models choose a simplicial approximation $g$ to $f$
after subdivision, as in
\cite[Theorem~2C.1, pp.~177--179]{Hatcher2002}.
The models and $g$ are semialgebraic.  Orient each model by its
identifying homeomorphism; these identifications preserve the
fundamental classes and the orientation forms used for degree and
pushforward.  The homotopy need only be continuous:
\Cref{thm:homotopy} identifies $f_*$ and $g_*$ on the continuous
theory, with the model identifications understood.  The constant
sheaf on the source model is an identity-map generator, so its
orientation form belongs to the semialgebraic category.

Fix a top dimensional target simplex.  In its interior choose
a small ball $B$ avoiding the images of lower dimensional source
simplices and of collapsed top dimensional simplices.  Each
source simplex mapping onto this target simplex contributes
one open sheet over $B$, mapped homeomorphically onto $B$.
These are all preimages of $B$, and there are finitely many.
The restricted map is proper, as the base change of the map
from the compact source.  On $B$ its direct image is the
constant vector space with one basis vector for each sheet.

The pushed pairing is diagonal in this basis.  Products between
different sheets vanish; on one sheet the trace of the
orientation pairing is $+1$ or $-1$ according as that
homeomorphism preserves or reverses the chosen orientations.
To see the sign directly, the trace for a homeomorphism
identifies its source dualizing complex with the target
dualizing complex; relative to the two orientation
trivializations, the induced map on local top homology is
exactly its orientation sign.  Proper duality and its trace
therefore give the same sign for the constant-line pairing.
Its signature on $B$ is the sum of these local signs.  For a sphere
source this is the degree $s$ by the local-degree formula
\cite[Proposition~2.30, pp.~135--136]{Hatcher2002}.
For the stated general source, fix $y\in B$ and write
$g^{-1}(y)=\{x_1,\ldots,x_t\}$.  Excision identifies
\[
 H_d(N,N\setminus g^{-1}(y);\Z)
 \simeq\bigoplus_{i=1}^t H_d(N,N\setminus\{x_i\};\Z).
\]
The fundamental class $[N]$ maps to the tuple of positive local
orientation generators.  The relative map induced by $g$ sends the
$i$th generator to $\epsilon_i$ times the orientation generator of
$H_d(S^d,S^d\setminus\{y\};\Z)$, where $\epsilon_i$ is the
orientation sign of that sheet.  Naturality of the absolute-to-relative
maps, together with $g_*[N]=s[S^d]$, gives
$s=\sum_i\epsilon_i$.  This also covers an empty fiber, with sum zero.
Only interiors of top dimensional simplices occur; these are open
balls, so no PL manifold structure is required of the triangulation.

This local calculation determines the rational coefficient
without assuming a globally defined generic signature.  In
fact, if the rational image of $g_*u_{F,d}$ (or of the pushed class
from $N$) is $\rho\overline{u}_{F,d}$, clear denominators and torsion to obtain a finite Witt
relation.  Choose the signature ball for that relation inside
the ball of sheets just used; this is possible since the
relation's common open set is dense.  Its signature says
$\rho=s$.  Apply the comparison to obtain the assertion in $E$.
Finally, $\overline{u}_d$ spans the rational group by \Cref{thm:sphere-lattice},
so the assertion about the entire group follows by linearity.
\end{proof}

\begin{proposition}[A contractible open set can be localized away]
\label{prop:contractible-quotient}
Let $d\geq3$ be odd. Let $X$ be any test space for $\T$, and let
$U\subset X$ be a
nonempty contractible open subset.  Then the quotient functor
induces an integral isomorphism
\begin{equation}
 \label{lat:contractible-isomorphism}
 E_d(X)\xrightarrow{\ \sim\ }
 \Witt_d\bigl(\T(X)/\T(U)\bigr).
\end{equation}
In particular this applies to compact test polyhedra and to
their manifold-with-boundary models.  It is natural for
continuous maps $g:X\to X'$ with $g(U)\subset U'$, whenever
$U'$ is also a nonempty contractible open subset.
\end{proposition}

\begin{proof}
Homotopy invariance and \Cref{lat:point-groups} give
$E_d(U)=E_d(\pt)=0$.  Let $i:U\hookrightarrow X$ and let
$p_Z:Z\to\pt$ denote projection.  The composite
\[
 E_{d-1}(U)\xrightarrow{i_*}E_{d-1}(X)
 \xrightarrow{(p_X)_*}E_{d-1}(\pt)
\]
is $(p_U)_*$, an isomorphism.  Thus $i_*$ is split injective,
with retraction $((p_U)_*)^{-1}(p_X)_*$.  If $d\equiv1\pmod4$,
the point group here is $\Z$; if $d\equiv3\pmod4$, it is zero.
The split-injectivity argument applies in either case.
All these maps exist
by compact-support pushforward, whether or not the space maps
are proper.  The relevant part of \Cref{thm:witt-localization} is
\[
 0\longrightarrow E_d(X)
 \longrightarrow\Witt_d\bigl(\T(X)/\T(U)\bigr)
 \longrightarrow E_{d-1}(U)\xrightarrow{i_*}E_{d-1}(X).
\]
Exactness and the last injectivity give the asserted
isomorphism over $\Z$.  The functor $g_*$ preserves the support
subcategory because $g(U)\subset U'$, so naturality follows
from the natural localization sequence.  No rationalization
or splitting of quotient idempotents is involved.
\end{proof}
